%\chapter{Hauptteil} %- 3 Seiten
%In einer theoretischen Hausarbeit wird von Ihnen zusätzlich verlangt,
%• Argumentationen und Positionen verschiedener Autoren gegeneinander abzuwägen, zu diskutieren und mit Literatur zu belegen,
%• den Lesenden den Aufbau und den roten Faden der Argumentation zu verdeutlichen.


%\section{Folgen von \glqq ohne Verschlüsselung\grqq{}}

%\chapter{CSA-Verordnung - der Gesetzesentwurf der EU}
\chapter{CSA-Verordnung}
\label{ch:CSA}

% ******************************
% Einleitungssatz & Begriffsdefinition
Aufbauend auf den im vorherigen Kapitel definierten Begriffen, soll im Zuge dieser Arbeit die aktuell diskutierte, ausstehende CSA-Verordnung in den Fokus gerückt werden.\\
Die CSA-Verordnung (Child Sexual Abuse) ist ein Vorschlag der Europäischen Kommission für eine Verordnung \glqq zur Festlegung von Vorschriften zur Prävention und Bekämpfung des sexuellen Missbrauchs von Kindern\grqq{} \cite{steinebach:2023:Medien} im Internet. Diese wurde im Jahr 2022? erstmals vorgeschlagen und dessen vorgesehene Umsetzung wird stark kritisiert. \cite{bundestag:2022, europa:2022, steinebach:2023:Medien}

% **************************************************************************************************
\section{Idee}   % Grundgedanke / Inhalt
\label{sec:CSA:Idee}

... Text ...
- auch in Echtzeit, nicht erst nur auf Verdacht
- gegen Darstellungen sexualisierter Gewalt an Kindern - gegen CSAM
- Grooming

% **************************************************************************************************
\section{Chatkontrolle - Vorschläge zur Umsetzung} % konkrete Vorschläge zur Umsetzung, wertfrei
%\subsection*{technische Umsetzung}
\label{sec:CSA:Vorschläge zur Umsetzung}
Die EU-Kommission hatte im Mai 2022 die Pläne für eine Chatkontrolle vorgelegt \cite{schulzki-haddouti:2023}
- technische Umsetzung der Risiko-Abschätzung
- Netzsperren

\glqq Chatkontrolle\grqq{} \cite{europa:2022}

Datenerhebung
o	Bekanntes Material
o	Unbekanntes Material (mittels Machine Learning)
o	Grooming (mehr Machine Learning) - gezielte Kontaktaufnahme Erwachsener mit Minderjährigen in Missbrauchsabsicht bezeichnet, indem stufenweise ihr Vertrauen erschlichen wird.
o	Metadatenanalyse

o	LIBE-Report: Keine Chatkontrolle, Schutz vor Verschlüsselung, Schutz von anonymer Kommunikation

% *********
-	Edps – European Data Protection Supervisor (europäischer Dateschutzbeauftragter)
  o	Aushebelung der sicheren Ende-zu-Ende Verschlüsselung
    	sehr einfach den Eingriff in die Privatsphäre auf den Endgeräten auf alle Bereiche auszuweiten
    	Chatkontrolle
  o	Altersverifikation: Ausweis oder Gesicht zeigen (beides nicht anonym)
    	Bedrohung des Rechts auf anonyme Kommunikation
  o	Netzsperren
    	wenn nicht möglich Missbrauchsmaterial zu löschen, soll der Zugriff drauf gesperrt werden (subjektive Meinung: untaugliche Methode)
  o	„digitaler Hausarrest“ – Alterscheck vor App-Installation mit Grooming Risiko (<16)
-	Grooming auf allen Platformen mit Kontaktaufnahmemöglichkeit möglich
-	Wäre so als würd man Kindern verbieten rauszugehen, es könnte sie ja jemand ansprechen



% **************************************************************************************************
\section{Kritik an der CSA-Verordnung}          % erstes Werten, aber aus Quellen - wer hat was dagegen
\label{sec:CSA:Kritik}

... Text ...

% ******************************************************
% Formulierungshilfen - Vergleich von Positionen
%
% Während [Wissenschaftler/-in] die Auffassung vertritt, dass […], ist [Wissenschaftler/-in] der Ansicht, dass […]
% Diese Frage wird kontrovers diskutiert. So stellte [Wissenschaftler/-in] fest, dass […], während [Wissenschaftler/-in] herausfand, dass […]
% [Wissenschaftler/-in] zeigte in seiner/ihrer Studie, dass […]. [Wissenschaftler/-in] kam hingegen zu dem Ergebnis, dass […]
